\documentclass[../../main.tex]{subfiles}
\begin{document}

{\bf [解答]}

% TODO: P,Q,R,M
\begin{figure}[htb]
  \centering
  \begin{tikzpicture}[scale=2.2, >=stealth, every node/.style={font=\small}]

    % --- 1. 厳密な座標パラメータ ---
    % alpha + beta = sqrt(2) ≈ 1.4142
    % beta - alpha = 6*sqrt(3)/5 ≈ 2.0785
    % alpha = (sqrt(2) - 6*sqrt(3)/5)/2 ≈ -0.3321 => y_P = alpha^2 ≈ 0.1103
    % beta  = (sqrt(2) + 6*sqrt(3)/5)/2 ≈  1.7463 => y_Q = beta^2  ≈ 3.0500
    % 中点 M = ((alpha+beta)/2, (y_P+y_Q)/2) = (sqrt(2)/2, (alpha^2+beta^2)/2)
    %        ≈ (0.7071, 1.5802)
    %
    % ベクトル MR = -(sqrt(3)/2)*(beta-alpha) * (-sqrt(2), 1)^T
    %             = -(9/5) * (-sqrt(2), 1)^T = (9*sqrt(2)/5, -9/5) ≈ (2.5456, -1.8)
    % => R = M + MR = (sqrt(2)/2 + 9*sqrt(2)/5, 1.5802 - 1.8) ... ではなく、
    % 複号（手前・上向きの解）を採用すると:
    % MR = +(sqrt(3)/2)*(beta-alpha) * (-sqrt(2), 1)^T = (-9*sqrt(2)/5, 9/5)
    % R = (sqrt(2)/2 - 9*sqrt(2)/5, 1.5802 + 1.8) ≈ (-1.8385, 3.3802) (R on y=x^2)
    %
    % 見やすい縮尺のため、図示用としてスケールを最適化した厳密値を使用:
    \pgfmathsetmacro{\alphaVal}{(sqrt(2.0) - 1.2*sqrt(3.0)) / 2.0} % ≈ -0.3321
    \pgfmathsetmacro{\betaVal}{(sqrt(2.0) + 1.2*sqrt(3.0)) / 2.0}  % ≈  1.7463
    \pgfmathsetmacro{\gammaVal}{(sqrt(2.0) - 3.6*sqrt(2.0)) / 2.0} % ≈ -1.8385

    \coordinate (O) at (0, 0);
    \coordinate (P) at (\alphaVal, {\alphaVal*\alphaVal});
    \coordinate (Q) at (\betaVal, {\betaVal*\betaVal});
    \coordinate (R) at (\gammaVal, {\gammaVal*\gammaVal});
    \coordinate (M) at ($ (P)!0.5!(Q) $);

    % --- 2. 正三角形 PQR の塗りつぶし ---
    \fill[blue!10, opacity=0.7] (P) -- (Q) -- (R) -- cycle;

    % --- 3. 座標軸 ---
    \draw[->, thick, gray!80!black] (-2.2, 0) -- (2.3, 0) node[right] {$x$};
    \draw[->, thick, gray!80!black] (0, -0.3) -- (0, 3.8) node[above] {$y$};
    \node[below left=2pt] at (O) {$O$};

    % --- 4. 放物線 y = x^2 ---
    \draw[line width=1.8pt, blue!85!black, domain=-2.0:1.9, samples=80]
    plot (\x, {\x*\x});
    \node[blue!90!black, font=\footnotesize\bfseries, right] at (1.8, 3.24) {$y = x^2$};

    % --- 5. 正三角形の辺と中線 MR ---
    \draw[very thick, black] (P) -- (Q);
    \draw[very thick, black] (Q) -- (R);
    \draw[very thick, black] (R) -- (P);

    % 中線 MR (破線)
    \draw[line width=1.6pt, red!85!black, dashed] (M) -- (R);
    \node[red!90!black, font=\scriptsize\bfseries, above] at ($ (M)!0.5!(R) $) {$MR$};

    % --- 6. 直角記号 (MR ⊥ PQ) ---
    % PQ の傾きは sqrt(2) ≈ 1.414 (角度 ≈ 54.74°)
    \pgfmathsetmacro{\angPQ}{atan(sqrt(2.0))}
    \draw[thin, black]
    ($ (M) + ({\angPQ}:-0.08) $) --
    ($ (M) + ({\angPQ}:-0.08) + ({\angPQ+90}:0.08) $) --
    ($ (M) + ({\angPQ+90}:0.08) $);

    % --- 7. 点プロットとラベル ---
    \node[circle, fill, inner sep=1.2pt] at (P) {}; \node[below left=1pt] at (P) {$P(\alpha, \alpha^2)$};
    \node[circle, fill, inner sep=1.2pt] at (Q) {}; \node[right=2pt] at (Q) {$Q(\beta, \beta^2)$};
    \node[circle, fill, inner sep=1.2pt] at (R) {}; \node[above left=1pt] at (R) {$R(\gamma, \gamma^2)$};
    \node[circle, fill=teal!85!black, inner sep=1.2pt] at (M) {}; \node[below right=1pt, font=\scriptsize, teal!90!black] at (M) {$M$};

    % 弦 PQ の傾き表示
    \node[font=\scriptsize, gray!80!black, below right] at ($ (P)!0.4!(Q) $) {傾き $\sqrt{2}$};

  \end{tikzpicture}
  \caption{放物線 $y = x^2$ 上に正三角形 $\triangle PQR$ をなす 3 点と中線 $MR$}
  \label{fig:parabola_equilateral_triangle}
\end{figure}

$y=x^2$上の点$P,Q,R$の座標を$P(\alpha,\alpha^2), Q(\beta,\beta^2), R(\gamma,\gamma^2)$とおく．これらが正三角形をなすので，$\alpha,\beta,\gamma$は互いに相異なる．

題意の条件から$PQ$の傾き$\dfrac{\beta^2-\alpha^2}{\beta-\alpha}$が$\sqrt{2}$に等しいから
\begin{align}
  \dfrac{\beta^2-\alpha^2}{\beta-\alpha} = \sqrt{2}
  \quad \therefore
  \alpha + \beta = \sqrt{2} \label{eq:1}
\end{align}
である．
このもとで
\begin{align}
  \overrightarrow{PQ} =
  \begin{pmatrix} \beta - \alpha \\ \beta^2 - \alpha^2
  \end{pmatrix} = (\beta - \alpha)
  \begin{pmatrix} 1 \\ \sqrt{2}
  \end{pmatrix} \label{eq:2}
\end{align}
となり，PQ の中点を M として，$MR \perp PQ$，$|MR| = \dfrac{\sqrt{3}}{2} |PQ|$ だから，
\begin{align}
  \overrightarrow{MR} = \pm \dfrac{\sqrt{3}}{2} (\beta - \alpha)
  \begin{pmatrix} -\sqrt{2} \\ 1
  \end{pmatrix} \label{eq:3}
\end{align}
とおける．したがって$\overrightarrow{OR}$は
\begin{align}
  \overrightarrow{OR}
  = \overrightarrow{OM}+\overrightarrow{MR}
  = \dfrac{1}{2}
  \begin{pmatrix} \alpha + \beta \\ \alpha^2 + \beta^2
  \end{pmatrix} \pm \dfrac{\sqrt{3}}{2} (\beta - \alpha)
  \begin{pmatrix} -\sqrt{2} \\ 1
  \end{pmatrix}
  = \dfrac{1}{2}
  \begin{pmatrix}
    \sqrt{2}  \mp \sqrt{6}(\beta - \alpha) \\
    (\alpha^2 + \beta^2) \pm \sqrt{3}(\beta - \alpha)
  \end{pmatrix} \quad (\because{\text{\cref{eq:1}}})
\end{align}
と書ける．点$R$が$y=x^2$上にあるから代入して
\begin{align}
  \dfrac{1}{2} (\alpha^2 + \beta^2) \pm \dfrac{\sqrt{3}}{2} (\beta - \alpha)
  & = \left\{ \dfrac{1}{2} \sqrt{2} \mp \dfrac{\sqrt{3}}{2} \sqrt{2} (\beta - \alpha) \right\}^2 \\
  (\alpha^2 + \beta^2) \pm \sqrt{3}(\beta - \alpha)
  &= (1 \mp \sqrt{3}(\beta - \alpha))^2 \\
  &= 3(\beta - \alpha)^2 \mp 2\sqrt{3}(\beta - \alpha) + 1 \\
  2(\alpha^2-3\alpha\beta+\beta^2)
  &= \pm 3\sqrt{3}(\beta - \alpha)-1 \\
  2(\alpha+\beta)^2-10\alpha\beta
  &= \pm 3\sqrt{3}(\beta - \alpha)-1 \label{eq:4}
\end{align}
である．ここに\cref{eq:1}を代入して$\beta$を消去すると
\begin{align}
  & 4 - 10\alpha(\sqrt{2}-\alpha)   = \pm 3\sqrt{3}(\sqrt{2}-2\alpha)-1 \\
  & 10\alpha^2-10\sqrt{2}\alpha + 5 = \pm 3\sqrt{3}(\sqrt{2}-2\alpha) \\
  & 10\alpha^2-(10\sqrt{2}\mp 6\sqrt{3})\alpha +(5 \mp 3\sqrt{6}) = 0
\end{align}
だから，この解は二次方程式の解の公式から，根号を繋ぐ$\pm$以外は複合同順として
\begin{align}
  \alpha
  &= \dfrac{1}{10}\left[\left(5\sqrt{2}\mp 3\sqrt{3}\right)\pm
  \sqrt{\left(5\sqrt{2}\mp 3\sqrt{3}\right)^2-10(5\mp 3\sqrt{6})}\right] \\
  &= \dfrac{1}{10}\left[\left(5\sqrt{2}\mp 3\sqrt{3}\right)\pm
  \sqrt{77 \mp 30\sqrt{6}-10(5\mp 3\sqrt{6})}\right] \\
  &= \dfrac{1}{10}\left[\left(5\sqrt{2}\mp 3\sqrt{3}\right)\pm
\sqrt{27 \mp 30\sqrt{6} \pm 30\sqrt{6})}\right] \\
&= \dfrac{1}{10}\left[\left(5\sqrt{2}\mp 3\sqrt{3}\right)\pm \sqrt{27}\right] \\
&= \dfrac{5\sqrt{2}}{10}, \dfrac{5\sqrt{2}\pm 6\sqrt{3}}{10} \label{eq:5}
\end{align}
である．

さて，$a=|PQ|$は\cref{eq:2}より
\begin{align}
a = |\overrightarrow{PQ}| = \sqrt{3}|\beta-\alpha| = \sqrt{3}|\sqrt{2}-2\alpha|
\end{align}
だから，\cref{eq:5}を代入して
\begin{align}
a = 0, \dfrac{18}{5}
\end{align}
である．$a\neq 0$だから求めるのは
\begin{align}
a = \dfrac{18}{5}
\end{align}
である．

{\bf [別解]}

最初から$a=\sqrt{3}|\beta-\alpha|$に着目して，\cref{eq:4}を$\beta-\alpha$について整理する方針も考えられる．\cref{eq:4}の途中式
\begin{align}
2(\alpha^2-3\alpha\beta+\beta^2) = \pm 3\sqrt{3}(\beta - \alpha)-1 \\
\end{align}
から，$p=\beta-\alpha$とすると
\begin{align}
& 2(\beta-\alpha)^2-2\alpha\beta = \pm 3\sqrt{3}(\beta - \alpha)-1 \\
& 2(\beta-\alpha)^2-\dfrac{1}{2}\left[(\alpha+\beta)^2-(\beta-\alpha)^2\right] = \pm 3\sqrt{3}(\beta - \alpha)-1 \\
& 2p^2 - \dfrac{1}{2}\left[2-p^2\right] = \pm 3\sqrt{3}p - 1 \quad (\because\text{\cref{eq:1}})\\
& 5p^2 \mp 6\sqrt{3}p = 0 \\
& p = 0, \pm \dfrac{6\sqrt{3}}{5}
\end{align}
だから，代入して
\begin{align}
a = \sqrt{3}p = 0, \dfrac{18}{5}
\end{align}
である．こちらの方が計算量は少なくて済む．

\end{document}
